%  B_technology.tex
\section{Technology Stack}
\label{app:technology}

\subsection{Core framework}

All experiments are implemented in \textbf{Python~3.11} using
\textbf{PyTorch}~\cite{pytorch2019} as the deep learning backend.
The workspace is managed by \texttt{uv}, a fast Python package installer and resolver that
replaces \texttt{pip}/\texttt{conda} for dependency management.
A single \texttt{uv sync --frozen} command restores the full pinned environment from
\texttt{uv.lock}.

\subsection{Model and training infrastructure}

\begin{itemize}[itemsep=2pt]
  \item \textbf{PyTorch}~\cite{pytorch2019}: the primary deep learning framework.
        All model code uses native \texttt{nn.Module} with standard training loops.
  \item \textbf{Torchvision}: provides the ResNet-50 backbone with
        ImageNet1K-V2 pretrained weights (\texttt{ResNet50\_\allowbreak Weights.DEFAULT}).
  \item \textbf{HuggingFace Accelerate}~\cite{huang2022accelerate}: provides
        device management, AMP (bfloat16) and gradient accumulation;
        the trainer integrates an \texttt{Accelerator} instance for
        device-agnostic mixed-precision training.
  \item \textbf{EMA (Exponential Moving Average)}: the embedding trainers maintain an
        EMA copy of model weights (decay~$= 0.999$) using a lightweight in-house
        implementation; the EMA model is used for all validation and test inference.
  \item \textbf{Early stopping}: implemented with configurable patience and minimum delta;
        the monitored metric is \recall for pair-based models and val/ROC-AUC for
        embedding models.
\end{itemize}

\subsection{Configuration management}

Experiments are configured via \textbf{Hydra}~\cite{yadan2019hydra}, a framework for
hierarchical configuration composition.  Each run composes one YAML file from each of
seven configuration groups:

\begin{center}
\ttfamily\small
\textup{data} / \textup{model} / \textup{loss} / \textup{trainer} /
\textup{evaluation} / \textup{preprocessing} / \textup{sweep}
\end{center}

Pre-composed experiment stacks are stored in \texttt{configs/experiment/} (10 files
covering all species$\times$background$\times$model combinations).  The global seed
(default 42) is propagated to all random operations.

\subsection{Experiment tracking}

\textbf{MLflow}~\cite{mlflow2018} is used for experiment tracking.  The
\texttt{bat\_tracking} façade logs:
\begin{itemize}[itemsep=2pt]
  \item All training hyperparameters (model, loss, trainer, data params from an
        allow-list to avoid Hydra internals).
  \item Per-epoch train/val metrics including loss, accuracy, F\textsubscript{1},
        ROC-AUC, TAR@FAR values, and \recall.
  \item Test-set metrics after training.
  \item Artefacts: the unified PDF report, saliency images, and permutation test outputs.
\end{itemize}
The local MLflow tracking server stores data in \texttt{mlruns/} at the repository root.
Run IDs (quoted throughout Section~\ref{sec:experiments}) uniquely identify each
experimental result.

\subsection{Hyperparameter optimisation}

\textbf{Optuna}~\cite{akiba2019optuna} provides hyperparameter search over validation
metrics only.  The search space for ArcFace covers:
margin $m \in [0.3,\,0.6]$, scale $s \in [32,\,96]$, and
$\mathrm{lr} \in \log\text{-unif}[10^{-3},\,3{\times}10^{-1}]$, over 20 trials.
The Siamese recall sweep covers Focal $\alpha\in[0.5,0.95]$, $\gamma\in[1,4]$,
and $\mathrm{lr}\in\log\text{-unif}[10^{-5},\,2{\times}10^{-4}]$, over 30 trials.
Trial state is persisted to \texttt{optuna\_studies.db} (SQLite).
The test split is never accessed during sweeps (enforced by a unit test mock).

\subsection{Preprocessing}

Face images are extracted from raw video by the \textbf{YOLO}-based
pipeline~\cite{redmon2016yolo, jocher2022yolov8}:
\begin{enumerate}[leftmargin=1.5em, itemsep=1pt]
  \item \texttt{YOLOPoseEstimator} (\texttt{face\_pose.pt}) detects facial keypoints.
  \item \texttt{FaceAligner} (NumPy/OpenCV) rotates the image so that the two primary
        keypoints are horizontal, standardising pose.
  \item \texttt{YOLOSegmenter} (\texttt{face\_seg.pt}) produces a binary face mask.
  \item Square crop around the mask bounding box with outer margin ratio 0.1; canvas
        padding ensures out-of-bounds crops remain valid.
  \item Optional \texttt{BackgroundGenerator} replaces or composites the background.
  \item Resize to target resolution and normalise to $[0, 1]$.
\end{enumerate}
Predictions are cached (512~MB in-memory LRU cache) to avoid redundant YOLO inference
during repeated runs.  Frame quality is scored by Laplacian variance (OpenCV) and stored
in the manifest.

\subsection{Interpretability and reporting}

Post-training interpretability is provided by three adapters in \texttt{bat\_interpretability}:
\begin{itemize}[itemsep=2pt]
  \item \textbf{\texttt{SiameseSaliencyAdapter}}: gradient-based input saliency for the
        pair network.
  \item \textbf{\texttt{GradCAMAdapter}}: class activation maps~\cite{he2016resnet} for
        embedding models, highlighting regions most discriminative for a given identity.
  \item \textbf{\texttt{EmbeddingProjectionAdapter}}: t-SNE and UMAP projections of
        test-set embeddings, coloured by identity, providing a 2-D snapshot of the
        embedding space geometry.
\end{itemize}
The unified PDF report is assembled by \textbf{reportlab} (pure-Python PDF generation)
and uploaded to MLflow as a run artefact.

\subsection{CLI surface}

Table~\ref{tab:cli} lists the complete \batcli command surface.

\begin{table}[h]
\centering
\caption{\batcli command surface.}
\label{tab:cli}
\setlength{\tabcolsep}{5pt}
\begin{tabular}{lp{8.5cm}}
\toprule
Command & Purpose \\
\midrule
\texttt{train}        & Compose config $\to$ train $\to$ test eval $\to$ interpretability $\to$ PDF $\to$ MLflow $\to$ permutation test $\to$ optional promote \\
\texttt{sweep}        & Optuna search on validation split \\
\texttt{evaluate}     & Re-run test evaluation against a saved checkpoint \\
\texttt{permutation-test} & Inference-mode permutation test on a configured run \\
\texttt{compare}      & Side-by-side MLflow run comparison HTML \\
\texttt{compare-champion} & Same report, left column auto-resolved via \texttt{bat\_tracking.get\_champion} \\
\texttt{promote}      & Promote a registered model if criterion beats incumbent \\
\texttt{build-manifest} & Walk a directory tree $\to$ split + write manifest CSV \\
\bottomrule
\end{tabular}
\end{table}

Invoked without arguments, \batcli presents an interactive picker (powered by
\texttt{InquirerPy}) from which any action can be selected.

\subsection{Quality assurance}

The sole quality gate is \texttt{pre-commit} (no CI), configured with:
\begin{itemize}[itemsep=1pt]
  \item \texttt{ruff} for linting and import sorting.
  \item \texttt{mypy} for incremental static type checking.
  \item Standard hooks (trailing whitespace, end-of-file, YAML/TOML syntax).
\end{itemize}
Unit tests are co-located with packages and run via \texttt{uv run pytest packages/}.
The test suite covers the identity-disjoint splitter (including the mock that confirms
the test split is never accessed by sweeps), manifest hash integrity, metric computations,
and named edge cases in the verification/identification protocols.
